\capitulo{2}{Objetivos del proyecto}

Este capítulo tiene como propósito definir el alcance y las metas fundamentales que guían el desarrollo de AgroSkopos. El proyecto surge de la necesidad de proporcionar una \textbf{interfaz de usuario avanzada} para la gestión de un vehículo autónomo en entornos agrícolas, priorizando la \textbf{precisión en la telemetría} y la \textbf{fiabilidad en el control remoto}.

\section{Objetivo General}

El objetivo principal de este Trabajo de Fin de Grado es el \textbf{diseño, implementación y validación} de una \textbf{plataforma web centralizada} para la monitorización, telemetría y control de un robot agrícola autónomo. La solución debe actuar como un \textbf{visor de precisión}, permitiendo al operador supervisar el estado del vehículo en tiempo real y gestionar misiones de trabajo de manera eficiente, garantizando la integridad operativa mediante un \textbf{entorno de simulación}.

\section{Objetivos Específicos}

Para alcanzar el objetivo general propuesto, se han definido una serie de metas específicas divididas en tres vertientes: funcionales, técnicas y de calidad y gestión.

\subsection{Objetivos Funcionales}

\begin{itemize}
    \item \textbf{Visualización Geográfica:} Integrar un sistema de mapas interactivos basado en Leaflet \cite{leaflet} que permita el seguimiento preciso de la ubicación del robot y la representación de rutas de misión.
    \item \textbf{Panel de Telemetría:} Desarrollar un panel dinámico que procese y muestre variables críticas de estado, tales como el nivel de batería, la velocidad de avance y las coordenadas GPS.
    \item \textbf{Control Remoto y Vídeo:} Implementar una interfaz de control manual para la operación directa del vehículo, integrando un flujo de vídeo en baja latencia para la supervisión visual.
    \item \textbf{Gestión de Misiones:} Crear un sistema de persistencia para el historial de misiones, permitiendo la consulta de datos históricos y el análisis de rendimiento de jornadas anteriores.
\end{itemize}

\subsection{Objetivos Técnicos}

\begin{itemize}
    \item \textbf{Arquitectura de Baja Latencia:} Configurar un sistema de comunicación cliente-servidor en tiempo real mediante canales bidireccionales (\textit{WebSockets}) para asegurar que la telemetría refleje de forma instantánea el estado del vehículo y permitir el control remoto sin retardos.
    \item \textbf{Gestión de Estado Global:} Utilizar Zustand \cite{zustand} para la administración del estado de la aplicación en el \textit{frontend}, garantizando una reactividad fluida sin las sobrecargas de rendimiento asociadas a otras librerías.
    \item \textbf{Simulación Cinemática:} Desarrollar un motor de simulación nativo en el \textit{backend} que emule el comportamiento físico y energético del robot, permitiendo validar la lógica de la aplicación sin dependencia inicial del \textit{hardware} físico.
    \item \textbf{Arquitectura Asíncrona y Colas de Mensajes:} Incorporar un sistema de colas basado en Redis \cite{redis} para procesar tareas pesadas en segundo plano (como el envío de correos electrónicos transaccionales y la gestión de \textit{tickets} de soporte), descongestionando el hilo principal del servidor.
    \item \textbf{Contenerización y Despliegue en la Nube:} Orquestar todo el ecosistema de \textit{software} (interfaz, lógica de servidor, proxy inverso y base de datos) mediante Docker \cite{docker} para asegurar la portabilidad, y desplegar la arquitectura completa en una máquina virtual alojada en Microsoft Azure \cite{azure}.
\end{itemize}

\subsection{Objetivos de Gestión y Calidad}

\begin{itemize}
    \item \textbf{Metodología Ágil (Scrum):} Implementar un marco de trabajo basado en Scrum \cite{scrum}, estructurado en iteraciones o \textit{sprints} de dos semanas. Este enfoque permite gestionar la variabilidad en la disponibilidad temporal y garantiza una validación continua de los requisitos.
    \item \textbf{Pruebas Automatizadas y de Carga:} Diseñar e implementar una estrategia integral de pruebas, abarcando pruebas unitarias (\textit{Vitest} \cite{vitest}), pruebas de integración (\textit{End-to-End} con \textit{Supertest} \cite{supertest}) y pruebas de estrés para evaluar el rendimiento bajo alta concurrencia mediante \textit{k6} \cite{k6}.
    \item \textbf{Integración Continua (CI/CD) y Calidad de Código:} Integrar plataformas de automatización como \textit{GitHub Actions} \cite{githubactions} para validar el código antes de su despliegue, y emplear análisis estático mediante SonarCloud \cite{sonarqube} para garantizar una deuda técnica nula y el cumplimiento estricto de \textit{Quality Gates}.
\end{itemize}
